\section{Discussion}

The results demonstrate a strong negative association between bullying frequency and Grade~4 reading achievement in both South Africa and Brazil. These associations remained statistically significant after controlling for learner, household and school characteristics. The findings therefore support the main hypothesis and are consistent with previous international and African evidence \citep{FRY20186,ozyildirim2024effect,anton_erxleben2016bullying}.

The coefficients followed different patterns as controls were introduced. In South Africa, the estimated bullying penalties became progressively smaller, whereas the Brazilian estimates remained comparatively stable. This suggests that observed learner, household and school characteristics account statistically for part of the South African association but explain less of the Brazilian association. It does not establish that the remaining association is causal, because unobserved differences may still influence both bullying and achievement. Nevertheless, the contrast cautions against treating upper-middle-income countries as educationally homogeneous.

Several mechanisms may help explain the negative association. Previous research also identifies absenteeism, reduced school engagement, poor concentration and feelings of insecurity as possible pathways linking bullying to achievement \citep{FRY20186,anton_erxleben2016bullying}. These mechanisms could cause bullied learners to fall behind academically, but the present analysis does not test which pathway is primarily responsible.

Although previous South African research finds that boys are generally more likely to experience bullying than girls \citep{juan2025risk}, the gender interactions in the present analysis were not statistically significant. These findings are not contradictory: gender may shape learners’ exposure to particular forms of bullying without changing the association between a given frequency of bullying and reading achievement.

In Brazil, higher SES appeared to reduce the weekly bullying penalty, although the monthly interaction was not significant. In South Africa, the monthly and weekly penalties became larger as SES increased. One possible explanation is that bullied high-SES learners fall further behind the comparatively high scores of non-bullied learners from similar backgrounds. This explanation remains tentative because the analysis cannot identify the mechanisms producing the interaction.

Taken together, the findings indicate that bullying should be treated as an educational concern as well as a learner well-being concern. This is consistent with the Sustainable Development Goal commitments to end violence against children and promote safe, non-violent and inclusive learning environments \citep{FRY20186}. Because the data are cross-sectional, the estimates should not be interpreted causally, and unobserved characteristics may influence both bullying exposure and reading achievement. Nevertheless, the common-sample and continuous-scale robustness checks produced similar patterns, reducing concerns that the main findings are solely artefacts of changing samples or the categorical bullying measure.
